%Time-stamp: <ter out 20 21:42:00 2026 yair@asterix>
\documentclass[aspectratio=169,11pt]{beamer}

\usepackage[brazilian]{babel}
\usepackage[utf8]{inputenc}
\usepackage[T1]{fontenc}
\usepackage{amsmath,amssymb,mathtools}
\usepackage{tikz}
\usetikzlibrary{calc,arrows.meta}

\usecolortheme{default}
\setbeamertemplate{navigation symbols}{}

% ---------- macros ----------
\newcommand{\define}[1]{\textbf{#1}}
\newcommand{\tri}[3]{%
  \draw (#1,#2) -- ++(-.45,-.9) -- ++(.9,0) -- cycle;
  \node at (#1,#2-.65) {#3};}

\input{C_macros}   % \E, \ProbE, \Ind, \paren

\newtheorem{teorema}{Teorema}
\newtheorem{proposicao}[teorema]{Proposição}
\theoremstyle{definition}

\title{\textit{Treaps} (resumido)}
\subtitle{Árvores de busca aleatorizadas}
\author{J.D.}
\institute{UFABC}
\date{\today}

\begin{document}

\begin{frame}
  \titlepage
\end{frame}

% =====================================================
\begin{frame}{O problema}
  \textbf{Dicionário}: manter $S\subset U$ sob busca, inserção e remoção.
  \pause
  \vspace{1ex}

  Árvores binárias de busca (ABB): custo proporcional à \emph{altura}.
  Inserções em ordem crescente $\Rightarrow$ altura $n-1$.
  \pause
  \vspace{1ex}

  \textbf{Ideia}: usar aleatoriedade para que a árvore se comporte como
  uma ABB \emph{construída em ordem aleatória}, sem depender da ordem
  das operações.
\end{frame}

% =====================================================
\begin{frame}{Definição}
  Uma \define{treap} é uma árvore binária em que cada nó $v$ tem
  $\mathrm{chave}(v)$ e $\mathrm{prioridade}(v)$ tais que
  \begin{itemize}
  \item é \textbf{ABB} nas chaves;
  \item é \textbf{heap de máximo} nas prioridades:
    $\mathrm{prioridade}(v)>$ prioridade de todo descendente próprio.
  \end{itemize}
  \vspace{1ex}
  %Supomos chaves distintas e prioridades distintas.
\end{frame}

% =====================================================
\begin{frame}{Exemplo}
  \centering
  \begin{tikzpicture}[scale=.85,
    no/.style={circle, draw, minimum size=1.0cm, inner sep=0pt},
    rot/.style={font=\itshape\small, inner sep=0pt},
    val/.style={font=\small, inner sep=0pt}]
    \node[no] (j) at (0,0)       {};
    \node[no] (i) at (-3,-1.7)   {};
    \node[no] (m) at (3,-1.7)    {};
    \node[no] (a) at (-4.5,-3.4) {};
    \node[no] (b) at (-1.5,-3.4) {};
    \node[no] (l) at (1.5,-3.4)  {};
    \draw (j) -- (i);
    \draw (j) -- (m);
    \draw (i) -- (a);
    \draw (i) -- (b);
    \draw (m) -- (l);
    \foreach \n/\letra/\valor in {j/j/31, i/b/26, m/m/27, a/a/13, b/i/25, l/\ell/14}{
      \draw (\n.45) -- (\n.-135);
      \node[rot] at ($(\n)+(-0.2,0.2)$) {$\letra$};
      \node[val] at ($(\n)+(0.2,-0.2)$) {$\valor$};
    }
  \end{tikzpicture}

  \vfill

  \small letras = chaves \quad números = prioridades \quad
  em-ordem: $a,b,i,j,\ell,m$
\end{frame}

% =====================================================
\begin{frame}{Unicidade}
  \begin{proposicao}
    Dado um conjunto de pares (chave, prioridade) com chaves distintas
    e prioridades distintas, existe uma \textbf{única} treap que o
    representa.
  \end{proposicao}
  \pause
  \textbf{Ideia}: a heap obriga a raiz a ser o nó de maior prioridade;
  a busca determina quais chaves vão para cada lado; aplique indução
  nas subárvores.
  \pause
  \medskip

  \textbf{Consequência}: a forma depende só de chaves e prioridades,
  \emph{não} da ordem de inserção.  Equivale a inserir os nós em ordem
  \emph{decrescente} de prioridade numa ABB vazia.
\end{frame}

% =====================================================
\begin{frame}{Rotações}
  Consertam a \textit{heap} sem estragar a ABB.

  \vfill

  \centering
  \begin{tikzpicture}[scale=.8, >={Stealth[length=3mm,width=3.5mm]},
      no/.style={draw,rectangle,minimum size=5mm,inner sep=1pt,fill=white}]
    \node[no] (y)  at (0,0)  {$y$};
    \node[no] (x)  at (-1,-1) {$x$};
    \draw (y) -- (x);
    \draw (y) -- (1,-1);   \tri{1}{-1}{$C$}
    \draw (x) -- (-1.6,-2); \tri{-1.6}{-2}{$A$}
    \draw (x) -- (-0.4,-2); \tri{-0.4}{-2}{$B$}
    \draw[->] (2.3,-0.8) to[bend left=40]
      node[above] {direita em $y$} (4.2,-0.8);
    \draw[->] (4.2,-1.3) to[bend left=40]
      node[below] {esquerda em $x$} (2.3,-1.3);
    \begin{scope}[xshift=6.5cm]
      \node[no] (x2) at (0,0)  {$x$};
      \node[no] (y2) at (1,-1) {$y$};
      \draw (x2) -- (y2);
      \draw (x2) -- (-0.6,-1); \tri{-0.6}{-1}{$A$}
      \draw (y2) -- (0.4,-2);  \tri{0.4}{-2}{$B$}
      \draw (y2) -- (1.8,-2);  \tri{1.8}{-2}{$C$}
    \end{scope}
  \end{tikzpicture}

  \medskip
  A ordem em-ordem $A,x,B,y,C$ é preservada; custo $O(1)$.
\end{frame}

% =====================================================
\begin{frame}{Operações}
  \begin{description}
  \item[Busca] como em ABB; custo $\propto$ altura.
  \pause
  \item[Inserção de $x$] insere como folha com prioridade
    \textbf{sorteada}; enquanto $x$ tem prioridade maior que a do pai,
    rotaciona $x$ para cima.
  \pause
  \item[Remoção de $x$] enquanto $x$ não é folha, rotaciona $x$ com o
    filho de \textbf{maior} prioridade; depois remove a folha.
  \end{description}
  \pause
  \medskip
  Número de rotações $\le$ profundidade (inserção) ou altura da
  subárvore (remoção) de $x$: todas as operações custam $O(\text{altura})$.
\end{frame}

% =====================================================
\begin{frame}{Aleatoriedade}
  \begin{itemize}
  \item Prioridade de cada nó: uniforme em $[0,1]$, sorteada na
    inserção, independente das demais. Com probabilidade $1$, as
    prioridades são distintas.
  \item \textbf{Adversário oblívio}: a sequência de operações é
    escolhida independentemente das prioridades.
  \end{itemize}
  \pause Pela \alert{unicidade}, após qualquer sequência oblívia de
  operações a árvore é \textbf{a treap do conjunto atual, com
    prioridades i.i.d.\ uniformes}.

  $\Rightarrow$ basta analisar uma treap com $n$ nós e prioridades
  i.i.d.\ uniformes.

  \pause

  Independentemente de como o conjunto de dados evoluiu, a árvore
  resultante é exatamente a mesma que seria obtida se criássemos uma
  Treap do zero com os elementos atuais.

  
\end{frame}

% =====================================================
\begin{frame}{Profundidade esperada}
  Seja $x_k$ o $k$-ésimo menor elemento e $X(j,k)=\{x_j,\dots,x_k\}$.
  \begin{proposicao}
    $x_j$ é ancestral próprio de $x_k$ $\iff$ $x_j$ tem a
    \textbf{maior prioridade} em $X(j,k)$.
  \end{proposicao}
  \pause
  Logo $\ProbE{x_j \text{ ancestral de } x_k}=\dfrac1{|k-j|+1}$ e
  \[
    \E{D_k}=\sum_{j<k}\frac1{k-j+1}+\sum_{j>k}\frac1{j-k+1}
    =H_k+H_{n-k+1}-2\le 2\ln n .
  \]
  \pause
  Portanto busca, inserção e remoção custam $O(\log n)$ \textbf{em
    esperança}.
\end{frame}

% =====================================================
\begin{frame}{Alta probabilidade}
  \begin{theorem}
    Numa treap com $n\ge2$ nós e prioridades i.i.d.\ uniformes, com
    probabilidade pelo menos $1-2n^{-7}$ a profundidade de todo nó é
    menor que $2\mathrm e^2\ln n<15\ln n$.
  \end{theorem}
  \pause
  \textbf{Ideia da prova}
  \begin{itemize}
  \item $D_k=L_k+R_k$ (ancestrais à esquerda e à direita); cada lado é
    soma de Bernoullis \emph{independentes} de parâmetros $1/i$.
  \item Chernoff: $\ProbE{L_k\ge \mathrm e^2\ln n}<n^{-8}$,
    e o mesmo para $R_k$.
  \item União sobre os $n$ nós.
  \end{itemize}
\end{frame}

% =====================================================
\begin{frame}{Resumo}
  \begin{itemize}
  \item Treap = ABB nas chaves + heap de máximo nas prioridades;
    \textbf{forma única}.
  \item Inserção e remoção por \textbf{rotações}, $O(1)$ cada.
  \item Prioridades aleatórias $\Rightarrow$ árvore equivale a ABB de
    ordem aleatória.
  \item $\E{D_k}\le2\ln n$.
  \item Altura $<15\ln n$ com probabilidade $\ge1-2n^{-7}$.
  \end{itemize}
\end{frame}

\end{document}

%%% Local Variables:
%%% mode: LaTeX
%%% TeX-master: t
%%% End:
